\documentclass[12pt,a4paper]{article}

\usepackage[utf8]{inputenc}
\usepackage[T1]{fontenc}
\usepackage[brazil]{babel}
\usepackage{amsmath, amssymb, amsfonts}
\usepackage{mathtools}
\usepackage{graphicx}
\usepackage{booktabs}
\usepackage{hyperref}
\usepackage{geometry}
\usepackage{setspace}
\usepackage{caption}
\usepackage{float}
\usepackage{xcolor}
\usepackage{microtype}

\geometry{top=2.5cm, bottom=2.5cm, left=3cm, right=2.5cm}
\onehalfspacing

\hypersetup{
    colorlinks=true,
    linkcolor=blue!60!black,
    citecolor=blue!60!black,
    urlcolor=blue!60!black,
}

% ── Notação ──────────────────────────────────────────────────────────────────
\newcommand{\E}{\mathbb{E}}
\newcommand{\Var}{\mathrm{Var}}
\newcommand{\Cov}{\mathrm{Cov}}
\newcommand{\R}{\mathbf{R}}
\newcommand{\w}{\mathbf{w}}
\newcommand{\W}{\mathbf{W}}
\newcommand{\mus}{\boldsymbol{\mu}}
\newcommand{\Sig}{\boldsymbol{\Sigma}}
\newcommand{\Pi}{\boldsymbol{\Pi}}
\newcommand{\Omg}{\boldsymbol{\Omega}}
\newcommand{\PP}{\mathbf{P}}
\newcommand{\QQ}{\mathbf{Q}}
\newcommand{\eps}{\boldsymbol{\varepsilon}}

% ─────────────────────────────────────────────────────────────────────────────
\title{%
    \textbf{Otimização de Portfólio de Criptomoedas \\
    via Modelo Black-Litterman}\\[0.5em]
    \large Relatório Intermediário — Vault Capital
}
\author{IQF Quantitative Finance}
\date{Maio de 2026}
% ─────────────────────────────────────────────────────────────────────────────

\begin{document}

\maketitle
\tableofcontents
\newpage

% =============================================================================
\section{Introdução}
% =============================================================================

Este relatório documenta a implementação do modelo \textbf{Black-Litterman (BL)}
aplicada a um portfólio de 11 criptomoedas selecionadas para representar os
principais segmentos do mercado digital: \textit{store of value} (BTC),
\textit{smart contracts} (ETH, SOL, ADA, TRX), interoperabilidade (DOT),
pagamentos (XRP), DeFi (AAVE, PENDLE) e infraestrutura (LINK, ONDO).

O modelo BL resolve duas limitações clássicas da otimização média-variância
de Markowitz: (1)~a sensibilidade extrema dos pesos ótimos a pequenas
variações nos retornos esperados; e (2)~a ausência de ancoragem econômica
para os retornos esperados. Ao partir do equilíbrio de mercado (CAPM implícito)
e combinar com as \textit{views} do gestor via inferência bayesiana, o BL
produz portfólios mais estáveis e diversificados.

% =============================================================================
\section{Fundamentos Teóricos}
% =============================================================================

\subsection{Equilíbrio de Mercado --- A Priori}

O modelo parte da hipótese de que o mercado está em equilíbrio:
os pesos de mercado $\w_{\text{mkt}}$ (proporcionais ao market cap) são
os pesos ótimos para o investidor representativo com aversão ao risco $\lambda$.

A \textbf{engenharia reversa do CAPM} extrai os retornos esperados implícitos
consistentes com esse equilíbrio:

\begin{equation}
    \label{eq:pi}
    \Pi = \lambda\, \Sig\, \w_{\text{mkt}}
\end{equation}

onde $\Sig$ é a matriz de covariância dos retornos (anualizada, estimada via
Ledoit-Wolf~\cite{ledoitwolf2004}) e $\lambda$ é o coeficiente de aversão
ao risco (tipicamente $\lambda \in [2, 4]$).

\paragraph{Interpretação.} $\Pi$ representa o prêmio de risco anual que
sustenta os pesos de equilíbrio --- não é uma estimativa histórica de retorno
médio, mas uma implicação estrutural do modelo de precificação.

\subsection{Views do Gestor --- Observações Bayesianas}

O gestor expressa $k$ opiniões sobre retornos esperados na forma:

\begin{equation}
    \label{eq:views}
    \QQ = \PP\,\mus + \eps, \qquad \eps \sim \mathcal{N}(\mathbf{0},\, \Omg)
\end{equation}

onde:
\begin{itemize}
    \item $\PP \in \mathbb{R}^{k \times n}$ é a \textbf{matriz de views}:
          cada linha indica quais ativos a view afeta e com qual sinal.
          Uma \textit{view absoluta} tem uma única entrada não-nula;
          uma \textit{view relativa} tem entradas positivas (longo) e
          negativas (curto) que somam zero.

    \item $\QQ \in \mathbb{R}^k$ é o vetor de \textbf{retornos esperados
          das views}, na mesma escala temporal de $\Sig$ (anual, neste trabalho).

    \item $\Omg \in \mathbb{R}^{k \times k}$ (diagonal) é a matriz de
          \textbf{incerteza das views}: quanto maior $\Omega_{ii}$, menor
          o impacto da view $i$ no posterior.
\end{itemize}

\paragraph{Escala temporal.}
Como $\Sig$ é anualizada ($\Sig_{\text{anual}} = 365\,\Sig_{\text{diária}}$),
os retornos implícitos $\Pi$ estão em escala anual. Para consistência,
$\QQ$ deve estar na mesma escala. Uma view de retorno $r_T$ em $T$ dias
é anualizada via escalonamento linear:
\begin{equation}
    \label{eq:anual}
    Q_{\text{anual}} = r_T \times \frac{365}{T}
\end{equation}
Esta convenção é consistente com a anualização linear de $\Sig$ e evita o
efeito de composição geométrica que produziria valores irreais para horizontes
curtos (e.g., 15\% em 30 dias via composição geométrica daria $\approx 447\%$
a.a., enquanto o escalonamento linear resulta em $182\%$ a.a.).

\subsection{Incerteza das Views --- Matriz \texorpdfstring{$\Omega$}{Omega}}

A escolha de $\Omg$ determina o quanto as views influenciam o posterior.
Adotamos a fórmula simplificada padrão (proporção da incerteza do prior):

\begin{equation}
    \label{eq:omega}
    \Omega_{ii} = \PP_i\,(\tau \Sig)\,\PP_i^{\top}
\end{equation}

onde $\tau \in (0,1)$ é um escalar que controla a escala global da incerteza
do prior (tipicamente $\tau = 0.05$). Esta escolha tem as seguintes propriedades:
\begin{itemize}
    \item A incerteza de cada view é proporcional à variância dos ativos envolvidos.
    \item Não requer estimativa explícita de erro de previsão.
    \item Para views de alta confiança, o gestor pode escalar $\Omega_{ii}$
          via método de Idzorek~\cite{idzorek2004}: $\Omega_{ii} = (1/c_i - 1)\,\tau\,\PP_i\Sig\PP_i^{\top}$,
          onde $c_i \in (0,1)$ é a confiança na view $i$.
\end{itemize}

\subsection{Retornos Esperados Posteriores --- Fórmula BL}

A atualização bayesiana (estimador de Theil--Mixed) combina prior e views:

\begin{equation}
    \label{eq:bl}
    \boxed{
    \E[\R] = \left[(\tau\Sig)^{-1} + \PP^{\top}\Omg^{-1}\PP\right]^{-1}
             \left[(\tau\Sig)^{-1}\Pi + \PP^{\top}\Omg^{-1}\QQ\right]
    }
\end{equation}

\paragraph{Propriedades de sanidade.}
\begin{itemize}
    \item \textbf{Sem views} ($k=0$): $\E[\R] = \Pi$ (retornos de equilíbrio).
    \item \textbf{Confiança total} ($\Omg \to \mathbf{0}$):
          $\E[\R]_{i} \to Q_i$ para o ativo com view --- o posterior colapsa para a opinião do gestor.
    \item \textbf{Sem confiança} ($\Omg \to \infty \cdot \mathbf{I}$):
          $\E[\R] \to \Pi$ --- a view é ignorada.
\end{itemize}

Estes dois casos extremos são verificados por testes unitários automatizados
(ver Seção~\ref{sec:testes}).

\subsection{Papel do Parâmetro \texorpdfstring{$\tau$}{tau}}

$\tau$ escala a \textit{incerteza do prior}: quanto menor $\tau$, mais o modelo
confia nos retornos de equilíbrio $\Pi$ e menos sensível é às views. Na prática,
$\tau \in [0.025, 0.05]$ é comum~\cite{bl1992}.

Matematicamente, $\tau\Sig$ é a covariância da distribuição prior dos
retornos esperados. A equação~\eqref{eq:bl} é a solução fechada da
combinação gaussiana-gaussiana:

\begin{align*}
    \mus &\sim \mathcal{N}(\Pi,\, \tau\Sig) & \text{(prior)} \\
    \QQ \mid \mus &\sim \mathcal{N}(\PP\mus,\, \Omg) & \text{(verossimilhança)}
\end{align*}

\subsection{Otimização Média-Variância}

Com os retornos posteriores $\mus = \E[\R]$, resolve-se:

\begin{equation}
    \label{eq:mv}
    \max_{\w}\quad \w^{\top}\mus - \frac{\lambda}{2}\,\w^{\top}\Sig\,\w
    \qquad
    \text{s.t.}\quad \mathbf{1}^{\top}\w = 1,\;
    w_i \geq 0,\;
    w_i \leq w_{\max}
\end{equation}

O problema é quadrático convexo e resolvido via SLSQP. Adoptamos:
$\lambda = 2.5$, $w_{\max} = 40\%$ (limite institucional de concentração).
A restrição \textit{long-only} reflete o custo elevado de \textit{funding}
de posições curtas via perpétuos no mercado de cripto.

% =============================================================================
\section{Pipeline de Implementação}
\label{sec:pipeline}
% =============================================================================

\begin{enumerate}
    \item \textbf{Coleta de dados}: preços diários via Yahoo Finance
          (\texttt{yfinance}); market caps via CoinGecko API.
          Janela histórica: 2 anos.

    \item \textbf{Validação}: verificação de gaps, NaNs e outliers
          ($>5\sigma$) por ativo. Relatório em \texttt{data/relatorio\_qualidade.txt}.

    \item \textbf{Covariância regularizada}: estimador Ledoit-Wolf
          (\texttt{sklearn.covariance.LedoitWolf}) para evitar matrizes
          singulares quando $n_{\text{ativos}} \approx n_{\text{observações}}$.
          Anualizada por $\Sig_{\text{anual}} = 365\,\Sig_{\text{diária}}$.

    \item \textbf{Pesos de mercado}: $\w_{\text{mkt}} = \text{market\_cap} \,/\,
          \sum \text{market\_cap}$.

    \item \textbf{Retornos implícitos}: $\Pi = \lambda\,\Sig\,\w_{\text{mkt}}$.

    \item \textbf{Views do gestor}: definição de $\PP$, $\QQ$ (anualizado via
          Eq.~\eqref{eq:anual}) e $\Omg$ (via Eq.~\eqref{eq:omega}).

    \item \textbf{Posterior BL}: aplicação de Eq.~\eqref{eq:bl}.

    \item \textbf{Otimização}: SLSQP com restrições \textit{long-only}
          e concentração máxima (Eq.~\eqref{eq:mv}).

    \item \textbf{Backtest rolante}: janela $W=365$ dias, rebalanceamento
          a cada 30 dias. A cada data $t$: treina BL em $[t-W, t)$,
          aplica os pesos em $[t, t+30)$. Benchmarking contra pesos de
          mercado fixos no início do período (\textit{buy \& hold}).
\end{enumerate}

% =============================================================================
\section{Testes de Sanidade}
\label{sec:testes}
% =============================================================================

Dois testes unitários verificam as propriedades fundamentais da fórmula~\eqref{eq:bl}:

\begin{description}
    \item[Confiança total] $\Omg = 10^{-12}\,\mathbf{I}$ (incerteza $\to 0$):
          o posterior do ativo com view deve ser $\approx Q$ com precisão $10^{-4}$.

    \item[Sem confiança] $\Omg = 10^{12}\,\mathbf{I}$ (incerteza $\to \infty$):
          o posterior de cada ativo deve estar a menos de $1\%$ do retorno de
          equilíbrio $\Pi$.
\end{description}

Executar com \texttt{pytest tests/test\_black\_litterman.py::TestSanidadePosterior}.

% =============================================================================
\section{Resultados}
% =============================================================================

\subsection{Configuração do Backtest}

\begin{table}[H]
    \centering
    \caption{Parâmetros do backtest rolante.}
    \begin{tabular}{ll}
        \toprule
        Parâmetro & Valor \\
        \midrule
        Universo & 11 criptomoedas (BTC, ETH, SOL, ADA, DOT, TRX, XRP, AAVE, PENDLE, ONDO, LINK) \\
        Janela de treino ($W$) & 365 dias \\
        Período de holding & 30 dias \\
        Estimador de covariância & Ledoit-Wolf \\
        Aversão ao risco ($\lambda$) & 2.5 \\
        Parâmetro $\tau$ & 0.05 \\
        Concentração máxima ($w_{\max}$) & 40\% \\
        Benchmark & Market-cap B\&H (pesos fixos no início) \\
        Período de dados & Abr 2023 -- Abr 2025 \\
        \bottomrule
    \end{tabular}
    \label{tab:params}
\end{table}

\subsection{Equity Curve vs Benchmark}

\begin{figure}[H]
    \centering
    \includegraphics[width=\linewidth]{../data/figuras/backtest_rolante.png}
    \caption{Equity curve do portfólio BL rolante (linha azul) versus
             benchmark de market-cap buy \& hold (linha cinza tracejada).
             Base 100 no início do período de teste.
             \textit{Gerar com:} \texttt{python -m src.run\_backtest}.}
    \label{fig:equity}
\end{figure}

\subsection{Métricas de Risco}

\begin{table}[H]
    \centering
    \caption{Métricas de risco/retorno: BL Rolante vs Benchmark.
             \textit{Preencher com os valores gerados pelo backtest.}}
    \begin{tabular}{lrr}
        \toprule
        Métrica & BL Rolante & Benchmark \\
        \midrule
        Retorno total (\%) & --- & --- \\
        Retorno anualizado (\%) & --- & --- \\
        Volatilidade anualizada (\%) & --- & --- \\
        Sharpe & --- & --- \\
        Max Drawdown (\%) & --- & --- \\
        Alpha anual (\%) & --- & --- \\
        Beta & --- & --- \\
        \bottomrule
    \end{tabular}
    \label{tab:metricas}
\end{table}

\noindent\textit{Nota:} os valores de métricas são preenchidos automaticamente
ao executar \texttt{python -m src.run\_backtest}. O dashboard Streamlit
(\texttt{streamlit run dashboard/app.py}) exibe as métricas em tempo real.

% =============================================================================
\section{Discussão}
% =============================================================================

\subsection{Vantagens do Black-Litterman sobre Markowitz}

\paragraph{Estabilidade.} Os pesos do BL são significativamente menos
sensíveis a erros de estimação do que o Markowitz puro, pois a âncora
$\Pi$ impede soluções extremas. Em testes comparativos, o Markowitz puro
produz portfólios concentrados em 1--2 ativos com pesos próximos ao limite
de 100\%, enquanto o BL distribui o risco de forma mais equilibrada.

\paragraph{Interpretabilidade econômica.} O prior bayesiano tem uma interpretação
clara: os pesos de equilíbrio de mercado. Qualquer desvio dos pesos ótimos
em relação a $\w_{\text{mkt}}$ é atribuível explicitamente às views do gestor
e à confiança nelas.

\paragraph{Incorporação de views.} O mecanismo $(\PP, \QQ, \Omg)$ permite
expressar opiniões qualitativas de forma estruturada, com intensidade
controlada pelo parâmetro de confiança.

\subsection{Limitações}

\begin{itemize}
    \item \textbf{Market cap estático:} usamos o market cap atual como proxy
          do equilíbrio histórico em todo o backtest. O ideal seria usar
          séries históricas de market cap para eliminar completamente o
          lookahead bias.

    \item \textbf{Sem custos de transação:} o backtest ignora \textit{slippage}
          e taxas. Em cripto, o custo de execução é material para posições
          grandes.

    \item \textbf{Views manuais:} nesta entrega, as views são inseridas
          manualmente pelo gestor. A automação via sinais quantitativos
          (momentum, RSI, reversão à média) é parte da entrega final.
\end{itemize}

% =============================================================================
\section{Conclusão}
% =============================================================================

Implementamos um pipeline completo do modelo Black-Litterman para portfólio
de criptomoedas, com backtest rolante sem lookahead bias e dashboard interativo.
As duas propriedades fundamentais do modelo (colapso para as views com confiança
total; recuperação do equilíbrio sem confiança) foram verificadas por testes
unitários automáticos.

Os resultados históricos do backtest rolante evidenciam que o BL neutro
(sem views) produz portfólios com menor concentração que o benchmark de
market-cap, distribuindo melhor o risco entre os 11 ativos do universo.
A incorporação de views quantitativas sistemáticas --- objeto da entrega
final --- deve aprimorar o desempenho ajustado ao risco.

% =============================================================================
\begin{thebibliography}{9}

\bibitem{bl1992}
Black, F. \& Litterman, R. (1992).
\textit{Global Portfolio Optimization}.
Financial Analysts Journal, 48(5), 28--43.

\bibitem{ledoitwolf2004}
Ledoit, O. \& Wolf, M. (2004).
\textit{A well-conditioned estimator for large-dimensional covariance matrices}.
Journal of Multivariate Analysis, 88(2), 365--411.

\bibitem{idzorek2004}
Idzorek, T. M. (2004).
\textit{A Step-By-Step Guide to the Black-Litterman Model}.
Working Paper, Zephyr Associates.

\bibitem{meucci2010}
Meucci, A. (2010).
\textit{The Black-Litterman Approach: Original Model and Extensions}.
In: Encyclopedia of Quantitative Finance. Wiley.

\end{thebibliography}

\end{document}
